\documentclass[11pt]{article}
\usepackage[margin=1in]{geometry}
\usepackage{amsmath,amssymb,amsthm,hyperref}
\newtheorem{proposition}{Proposition}
\title{Arithmetic elimination and convolution in the C38 join}
\author{Project research note}
\date{17 November 2026}
\begin{document}
\maketitle

We use the project's general-prime C38 chord equations, symmetric
abscissa hashes and ordinary affine addition guards. The local elimination
below also works with a weighted source. For a whole undeformed chain the
guards admit an additive correction to convolution whose rank grows at
most linearly in the total slot-set size. The remaining cost is access to
coefficients of products in the curve group algebra.

\section{An exact join of the two chord equations}
Work over $\mathbb F_p$ for $p>3$, on a nonsingular curve
$y^2=f(x)=x^3+ax+c$. Write $u=X_{j-1},v=x_j,w=X_j$ and let $s$ be
the slope. Fix $u,w$ and the ordinate $y$ of the anchored point, at $w$
for an even step and at $u$ for an odd step. Let $\epsilon=1$ for an
even step and $-1$ for an odd step. The C38 equation $E_j=0$ sets
$v=s^2-u-w$. Substitution in the corresponding $D_j=0$ gives exactly
\[
  (w-u)s^2+2ys+\epsilon(u^2+uw+w^2+a)=0 .
\]
With $g=w-u$ and $h=u^2+uw+w^2+a$ we have $gh=f(w)-f(u)$.
Consequently $y^2-\epsilon gh$ equals $f$ at the other abscissa.
When $g\ne0$ the slopes are
\[
 s=\big(-y\pm\sqrt{y^2-\epsilon gh}\big)/g .
\]
Use one root when the radicand is zero and none when it is a nonsquare.
When $g=0$ and $y\ne0$, the unique candidate is
$s=-\epsilon h/(2y)$. When $g=y=0$ the anchored point is a root of
$f$ and nonsingularity gives $h=f'(u)\ne0$, so there is no solution.
These cases require an on-curve anchor. A synthetic off-curve boundary
with $g=y=h=0$ admits every slope in the two equations alone.

For each candidate enforce $v\ne u$ and the actual slot hash on the
canonical representative of $v$. Any additional link test or weight
depending on $(v,s)$ can be evaluated at the surviving pairs.
Thus two weight evaluations suffice, in addition to expected
$\operatorname{poly}(\log p)$ bit work for the field arithmetic. This
price assumes access to the specified boundary values and separately
charges evaluation of the weight. It does not sum over the boundary.
For completeness, a square-root table is unnecessary at large $p$.
Test for a nonzero square $d$ by Euler's criterion, then choose $t$
until $e^2=t^2-d$ defines a quadratic extension. There are $(p-1)/2$
suitable $t$, by counting $(t-z)(t+z)=d$. Frobenius sends $e$ to
$-e$, hence $(t+e)^{(p+1)/2}$ squares to $d$ and is one of its two
roots in $\mathbb F_p$. Binary powering gives the stated cost.

\paragraph{A deformed, source-weighted packet.}
Contract each such fixed-boundary chord slice first and encode its two
remaining field words in binary. If both the source and the target of
a packet are such slices, their analytically computed supports $I,J$
each have at most two elements. Let $W_{xy}(T)$ be the native packet
matrix, and let $w_X,v_Y$ be the scalar source and target weights.
For live positions the joined value is
\[
 \sum_{X\in I,Y\in J}w_Xv_Y
 \prod_x
 \begin{cases}
 u_x(X),&Y_x=1,\\
 1-u_x(X),&Y_x=0,
 \end{cases}
 \qquad u_x(X)=\sum_{z} W_{xz}(T)X_z .
\]
This follows by substituting the actual binary selectors into the
deformation. It evaluates at most four products, charges
$O(L^2)$ ring operations at width $L$ once the jet kernel is built, and
includes repeated powers introduced by the mixing. The native kernel
can itself be constructed in polynomial time in explicit $L$, precision
and jet length by its finite products. This test merges equations
before forming the packet. The original nine-group recursion can split
their Fourier ports across groups, so this identity alone supplies no
closure result for its general child descriptors.

\section{The whole guarded chain}
Let $G=E(\mathbb F_p)$ with identity $O$, let $F_j$ be the actual signed
slot set, let $s_j=|F_j|$, and write $h_j$ for its indicator.
It satisfies $h_j(O)=0$ and $F_j=-F_j$. Write $d$ for the vector of
incoming path counts supported on affine points. At step $j$ first form
the ordinary group convolution $g=h_j*d$. Then the exact outgoing
vector is
\begin{equation}
 T_j d=g-g(O)\delta_O-
 \sum_{\substack{A\in F_j\\2A\ne O}}d(A)\delta_{2A},
 \qquad g(O)=\sum_{A\in F_j} d(-A).
 \label{guard}
\end{equation}
Indeed, an excluded step from $U$ with addend $A$ has equal abscissas,
so $A=-U$ or $A=U$. The first case lands at $O$. The second lands at
$2A$ and has precisely the incoming multiplicity $d(A)$. Excluding
$2A=O$ from the last sum counts a two-torsion exception once.
Repeated doubling images are added with their full multiplicity.
The derivation uses the original C38 bijection from valid chains to
canonical arithmetic assignments. It applies on the semantic point
boundary before Fourier-word deformation.

Extend the displayed linear operator to arbitrary vectors by the same
formula. Genuine input vectors to a step have zero at $O$. The start
of the full relation chain is $d_1=h_1$. Write $C_j$ for ordinary
convolution by $h_j$. Since the image of $T_j-C_j$ lies in the span
of the at most $1+s_j$ indicated point masses, telescoping gives
\[
 \operatorname{rank}(T_k\cdots T_i-C_k\cdots C_i)
 \le \sum_{j=i}^k(1+s_j).
\]
This follows over every field. At each telescoping summand, multiplication
on the left and right preserves the upper bound on its rank.
The ordinary part is a single convolution by $h_i*\cdots*h_k$.

\paragraph{Entry access and construction cost.}
Call its coefficient at $P$ $H_{i:k}(P)$, with the empty interval
equal to $\delta_O$. For each $i\ge2$ save the sparse correction $e_i$
from (\ref{guard}) using the actual preceding counts. The exact count is
\[
 d_j(W)=H_{1:j}(W)+
 \sum_{i=2}^j\ \sum_{P\in\operatorname{supp} e_i}
       e_i(P)H_{i+1:j}(W-P).
\]
To construct $e_i$, use this formula for $d_{i-1}$ at the $s_i$
members of $F_i$. Symmetry supplies all $d_{i-1}(-A)$ as well.
Let $D=1+\sum_i(1+s_i)$. Building the correction lists takes $O(D^2)$
interval-coefficient queries and arithmetic operations, then one final
endpoint needs $O(D)$ such queries and operations. Apart from storage
for slot sets, the coefficient procedure and answers, $O(D)$ locations
and integer weights suffice. Counts in an interval are bounded by the
product of its slot sizes. The bit lengths of intermediate sums here
are polynomial in $m$ and $b$ because the corrections are counts of
excluded transitions and there are at most $D$ terms.
Constructing slot sets also costs time, as in the original budget.

For two slots a coefficient can be obtained by scanning the smaller
set and testing membership of the complementary point in the other
hash. At longer intervals the formula still needs coefficients of an
ordinary product of many independent slot indicators. The verifier
builds their full finite tables and counts the group additions used.
On the order-19 five-slot fixture this costs 724 additions before the
607 coefficient lookups for correction setup and all target counts.
Dense tables on a group of size comparable to $2^b$ have exponential
size in $b$. The formulas supply no subexponential entry procedure for
the long intervals. They compute counts, and the full application would
also have to pay for witness selection, training, and extraction.

\section{Why a small sum of separated boundary functions fails}
Consider the prime-order case $|G|=r$ and a nonempty slot of size
$s$. This is an admissible subclass, and the following bounds apply to
each such finite group. On the whole group the unguarded convolution
matrix $C$ is invertible over $\mathbb Q$. Indeed, identify the group
abstractly with $\mathbb Z/r$ for the proof. The eigenvalues are the
evaluations at $r$th roots of unity of the degree-at-most-$(r-1)$
indicator polynomial. At one its value is $s$. At any other root,
vanishing would imply divisibility by $1+X+\cdots+X^{r-1}$, irreducible
by applying Eisenstein after translating $X$ by one. The indicator is
nonconstant since $O$ is excluded. This proof uses no computed discrete
logarithms. The character diagonalization of convolution is standard.
See \href{https://www2.math.upenn.edu/~gressman/analysis/14-discrete.html}
{P.\ Gressman, Fourier analysis on finite abelian groups}.

Keep the $r-1-s$ rows $U$ with $U\ne O$ and $U\notin F_j$ in the
matrix indexed by incoming and outgoing points. They are unchanged by
the C38 guard and have zero in column $O$. They were independent in
$C$ and remain so on the affine columns. Therefore
\[
 \operatorname{rank}_{\mathbb Q} K_j \ge r-1-s_j,\qquad
 \operatorname{rank}_{\mathbb Q}(K_i\cdots K_k)
 \ge r-1-\sum_{j=i}^k s_j ,
\]
where $K_j$ is the genuine affine-to-affine transition. The product
bound follows by adding nullities on the common $(r-1)$-dimensional
space. There is also a one-step bound in any characteristic. Each
retained row has support $s$. Choosing a single fixed $A\in F_j$ shows
that the union of these row supports contains at least $r-1-s$ distinct
columns $U+A$. A basis of retained rows must cover that union.
Consequently rank is at least $\lceil(r-1-s)/s\rceil$, including
over $\mathbb F_2$.

A representation $\sum_{\ell=1}^q a_\ell(U)b_\ell(W)$ must therefore
use at least the indicated rank in separated terms. For sparse slot sets
the rational bound can be nearly the full group order even after
composing steps. Convolution itself is a short description of such a
high-rank matrix. The result leaves open a nonlinear or implicit
algorithm for selected interval coefficients.

\paragraph{A cost bound for two explicit implementations of the remaining sum.}
Use the original targetwise coverage lemma at its prescribed hash
threshold on a prime-order curve with $r\ge30m$. Under its independent
seed law every nonidentity target has a guarded tuple with probability
greater than $2/5$. Let $Z$ be the number of such reachable targets.
Linearity gives $\mathbb E Z>2(r-1)/5$. Since $0\le Z\le r-1$,
\[
 \Pr\{Z\ge(r-1)/5\}\ge 1/4.
\]
Every guarded tuple is also an unguarded tuple, with positive integral
weight. On this event the support of $H_{1:m}$ contains at least
$(r-1)/5$ points. For every split $1\le j<m$,
\[
 |\operatorname{supp}H_{1:j}|\,|\operatorname{supp}H_{j+1:m}|
 \ \ge\ |\operatorname{supp}H_{1:m}|\ \ge\ (r-1)/5
\]
because full sums are sums of two partial sums. Constructing both
explicit partial-support dictionaries therefore needs at least
$\sqrt{(r-1)/5}$ distinct entries even if the split is chosen from the
data. In the character basis the full convolution has all $r$ Fourier
coefficients nonzero whenever all slots are nonempty, by the same
prime-order argument applied to each factor. Hence sparse tables in
either of these two bases cannot deliver a $2^{o(b)}$ expected build
cost on this subclass, where $r$ is comparable to $2^b$. This
consequence uses the project's stated coverage conditions. It concerns
algorithms materializing these supports, leaving an on-demand algorithm
for a single coefficient open.

\paragraph{Validation and conclusion.}
Run \texttt{python -B research/c38\_arithmetic\_join/verify\_join.py}.
The audit compares 810 fixed-boundary slices to both vendor equation
orientations across 116,442 raw pairs. It checks a mixed packet with eight
live bits in sixteen slots, two analytic chord supports, nonuniform weights,
and five jet coefficients modulo $2^{60}$. The four weighted support
pairs match all 270 admissible normalized exponent vectors. Dropping
repeated exponents changes the answer. Five finite chain fixtures include
empty slots, two-torsion, and curves of orders 14, 18, 19 and 31. Their
guarded counts agree at every target with literal ordered tuples. Every
accepted tuple is also checked against the original integer equations
including the shifted quotients. Separately, an order-223 example over
$\mathbb F_{211}$ meets the project's actual coverage threshold for four
slots. Both two-slot sumsets fill all 223 group elements and every one
of the 222 affine targets has a guarded path. The probability bound
above follows from the coverage lemma independently of this example.
These diagnostics do not run the full Fourier graph.
The conditioning cost and the missing fast general interval coefficient
procedure prevent any improvement here to the project's conditional
exponent or its leading constant.
\end{document}
